% Рисунок с доски: произвольный элемент α с нагрузкой S_α и деформацией e_α.
\begin{figure}[htbp]
\centering
\begin{tikzpicture}[>={Latex[length=4pt]}]
  \draw[thick] (-0.8,-0.5) rectangle (0.8,0.5);
  \node at (0,0) {$\alpha$};
  \draw[->, thick] (0,1.3) -- (0,0.55);
  \draw[->, thick] (0,-0.55) -- (0,-1.3);
  \node[anchor=east, font=\footnotesize] at (-0.1,1.1) {$\Sdev_\alpha$};
  \node[anchor=east, font=\footnotesize] at (-0.1,-1.1) {$\Sdev_\alpha$};
  \node[anchor=west, font=\footnotesize] at (0.1,1.1) {$\dev_\alpha$};

  \node[anchor=west, align=left, font=\footnotesize] at (2.0,0.5)
    {$\alpha = 1, 2, 3, \dots, N$};
  \node[anchor=west, align=left, font=\footnotesize] at (2.0,-0.4)
    {$\alpha$ может быть:};
  \pic at (6.0,-0.4) {spring};
  \pic at (7.9,-0.4) {damper};
  \pic at (9.8,-0.4) {stvenant};
\end{tikzpicture}
\caption{Элемент реологической модели. Комбинируя элементы, собираем новый
материал.}
\label{fig:rheo-block}
\end{figure}
